\section{相关工作}
\label{sec:related-work}

\subsection{表征工程与激活引导}

表征工程把群体层面的内部状态作为读出和控制对象，连接了线性探针、概念方向与模型干预\cite{zou2025representation}。Park 等从反事实语义形式化线性表征，并说明几何操作依赖所选内积\cite{park2024linear}。对比激活加法则从正负样本对的残差差异估计控制向量，在推理阶段将向量加入隐藏状态\cite{turner2023activation}；Rimsky 等在 Llama 2 上系统研究了这一路径\cite{rimsky2024steering}。这些方法证明方向性干预可以成为行为控制工具，却不保证一个方向能完整覆盖复杂安全行为。

激活加法与方向消融共享向量表示，却执行相反类型的操作。前者以带符号系数加入控制向量，允许连续调节目标行为；后者把向量张成的分量从激活或权重中删除。二者都依赖方向估计和干预层选择，也都可能影响非目标行为。本文把它们归入同一表征控制背景，但只将拒答投影方法作为直接结构基线。

\subsection{拒答的单方向与多维结构}

拒答方向研究提供了从干预到权重编辑的直接基线。Arditi 等报告单一方向在 13 个开源聊天模型上介导拒答，并将残差正交化折叠为紧凑权重修改\cite{arditi2024refusal}。安全微调位移分析与概念锥研究从不同数据锚点报告多维结构\cite{pan2025hidden,wollschlager2025geometry}；后者还提出表征独立性以区分几何正交与干预独立。Piras 等进一步使用 SOM 原型构造相关但不强制正交的方向池，并搜索有序投影组合\cite{piras2026som}。本文不重做这些论文的模型实验，而是把它们作为矩阵源码形式化的结构参照。

多方向证据并未形成唯一的拒答几何定义。概念锥工作从梯度信号和干预相互作用出发，SOM-MD 则从有害提示的局部原型出发；一个强调机制独立，一个允许方向保持较高相关性。二者都反对把欧氏正交直接当作完整因果解释，但其方向集合不可互换。矩阵法则改变了决策对象，因此本文只讨论结构差异，不把“无显式方向”解释成对所有多方向机制的吸收。

\subsection{局部原型、黑盒搜索与评测}

SOM 用规则网格上的原型近似高维数据，并保持原型间的邻接关系\cite{kohonen2013essentials}。贝叶斯优化把昂贵评估视为黑盒，TPE 用条件密度建模建议新的离散或连续配置\cite{snoek2012practical,bergstra2011algorithms}；Optuna 提供了该搜索的软件框架\cite{akiba2019optuna}。在拒答评测中，HarmBench 标准化自动红队与判别器流程，SORRY-Bench 则强调细粒度风险类别、语言和格式变化\cite{mazeika2024harmbench,xie2025sorrybench}。这些基准有助于定义统一复现实验，但不能替代对训练、搜索和测试样本关系的公开说明。

黑盒搜索的试验数只有在单次试验成本明确时才可比较。SOM-MD 的一次试验需要用一个有序方向组生成回答并调用行为判别器；矩阵法的一次外层试验还包含多个模块的 L-BFGS 优化。TPE 在两种流程中只是建议器，不会使二者自动获得相同的计算预算或统计独立性。评测论文提供了判别器和提示规范，却不能修复上游论文未公开的数据划分。

\subsection{本文位置}

本文位于方法论文与软件审计之间。我们不提出新的方向估计器，也不把源码原型的公开演示当作实证贡献。本文从固定实现分别重建上游硬近邻和本地 point-v1 软邻域目标，再以共同坐标比较显式方向投影与直接映射，并审计凸性、秩语义和实现一致性。新增的单次归档搜索呈现效果门槛失败，属于本地观察而非同行评议文献结果。这个定位把文献、实现、数学、归档观察与待验证假设分开。

现有工作留下的关键缺口不是再为一种参数化增加孤立效果数字，而是建立可复核的统一协议。统一协议需要同时控制模型、提示划分、干预位置、解码、判别器和总计算预算，并报告拒答抑制与无害能力保持。本文的比较表为这类实验定义对象；当前稿件完成形式化、证据审计及 point-v1 失败搜索分析，尚未完成统一协议的跨方法实验；trajectory-v2 的源码控制也尚未得到指定归档中的效果验证。
